% GENERATED FILE. Edit canonical lessons, then run npm run generate:books.
% canonical-source-hash: fnv1a64:f4699e4ef37155d2
% canonical-lessons: TA-C392-ulal, TA-C392-lancam, TA-C392-velippataiyaka, TA-C392-maraimukamaka, TA-C392-kankani, TA-R392-review-386-392

\chapter{Clean Government}
\label{ch:ta-clean-government}
\hlchaptermodality{392}

\begin{chapteropening}
\textbf{By the end of this chapter:} \emph{I can talk about corruption, bribes and keeping watch.}

Complete the last lesson of chapter 392: I can talk about corruption, bribes and keeping watch.
\end{chapteropening}

\section[\texorpdfstring{\ta{ஊழல்} (ūḻal)}{ஊழல் (ūḻal)}]{\ta{ஊழல்} (ūḻal) — corruption}

\label{lesson:TA-C392-ulal}



\begin{quote}
Before the new one: say the Tamil for to distribute, then the Tamil for to fix.
\end{quote}

\subsection*{You'll want to know: \ta{ஊழல்}}
\textbf{\ta{ஊழல்}} — \emph{ūḻal} — \textquotedblleft{}corruption\textquotedblright{}.

\begin{quote}
\ta{ஊழல்} \ta{இருந்தால்} \ta{மக்களுக்கு} \ta{அரசின்} \ta{மீது} \ta{நம்பிக்கை} \ta{போய்விடும்}.
\end{quote}

\begin{quote}
\emph{ūḻal iruntāl makkaḷukku araciṉ mītu nampikkai pōyviṭum.}
\end{quote}

\begin{quote}
If there is corruption, people lose faith in the government.
\end{quote}

\ta{ஊழல்} \ta{செய்}, to be corrupt; \ta{ஊழல்வாதி}, a corrupt person.

\begin{cousinweb}
\textbf{\ta{ஊழல்}} (\emph{ūḻal}) is native Dravidian: DEDR 756 lists Tamil \emph{ūḻal}, \textquotedblleft{}dirt, mire, that which is putrid, decayed, spoiled\textquotedblright{}: corruption is rot.
\end{cousinweb}

\begin{grammarlens}[title={What you've built}]
One more word for saying what happens if.
\end{grammarlens}

\subsection*{Your turn}
\begin{itemize}
\raggedright
  \item \emph{Say it:} \emph{ūḻal}
  \item \emph{Say it:} \emph{ūḻal}, once more
  \item \emph{Say it:} \emph{ūḻal iruntāl makkaḷukku araciṉ mītu nampikkai pōyviṭum.} — If there is corruption, people lose faith in the government.
  \item \emph{Recall:} say the Tamil for to distribute, then the Tamil for to fix, then say \emph{ūḻal} again
\end{itemize}

\begin{tcolorbox}[breakable,colback=teal!4,colframe=teal!35!black,title={Before you move on}]
What does \ta{ஊழல்} mean? (\textquotedblleft{}Corruption\textquotedblright{}.) Say it once more.
\end{tcolorbox}
\section[\texorpdfstring{\ta{லஞ்சம்} (lañcam)}{லஞ்சம் (lañcam)}]{\ta{லஞ்சம்} (lañcam) — a bribe}

\label{lesson:TA-C392-lancam}



\begin{quote}
Before the new one: say the Tamil for to fix, then the Tamil for corruption.
\end{quote}

\subsection*{You'll want to know: \ta{லஞ்சம்}}
\textbf{\ta{லஞ்சம்}} — \emph{lañcam} — \textquotedblleft{}a bribe\textquotedblright{}.

\begin{quote}
\ta{அவர்} \ta{லஞ்சம்} \ta{வாங்கினால்} \ta{வேலையை} \ta{இழப்பார்}.
\end{quote}

\begin{quote}
\emph{avar lañcam vāṅkiṉāl vēlaiyai iḻappār.}
\end{quote}

\begin{quote}
If he takes a bribe, he will lose his job.
\end{quote}

\ta{லஞ்சம்} \ta{வாங்கு}, to take a bribe; \ta{லஞ்சம்} \ta{கொடு}, to give one.

\begin{cousinweb}
\textbf{\ta{லஞ்சம்}} (\emph{lañcam}) is borrowed from Sanskrit \emph{lañcā}, \textquotedblleft{}a present, a bribe\textquotedblright{} (Monier-Williams), with the Tamil noun ending -\ta{அம்}.
\end{cousinweb}

\begin{grammarlens}[title={What you've built}]
One more word for saying what happens if.
\end{grammarlens}

\subsection*{Your turn}
\begin{itemize}
\raggedright
  \item \emph{Say it:} \emph{lañcam}
  \item \emph{Say it:} \emph{lañcam}, once more
  \item \emph{Say it:} \emph{avar lañcam vāṅkiṉāl vēlaiyai iḻappār.} — If he takes a bribe, he will lose his job.
  \item \emph{Recall:} say the Tamil for to fix, then the Tamil for corruption, then say \emph{lañcam} again
\end{itemize}

\begin{tcolorbox}[breakable,colback=teal!4,colframe=teal!35!black,title={Before you move on}]
What does \ta{லஞ்சம்} mean? (\textquotedblleft{}A bribe\textquotedblright{}.) Say it once more.
\end{tcolorbox}
\section[\texorpdfstring{\ta{வெளிப்படையாக} (veḷippaṭaiyāka)}{வெளிப்படையாக (veḷippaṭaiyāka)}]{\ta{வெளிப்படையாக} (veḷippaṭaiyāka) — frankly; openly}

\label{lesson:TA-C392-velippataiyaka}



\begin{quote}
Before the new one: say the Tamil for corruption, then the Tamil for a bribe.
\end{quote}

\subsection*{You'll want to know: \ta{வெளிப்படையாக}}
\textbf{\ta{வெளிப்படையாக}} — \emph{veḷippaṭaiyāka} — \textquotedblleft{}frankly; openly\textquotedblright{}.

\begin{quote}
\ta{பிரச்சினை} \ta{இருந்தால்} \ta{அவர்} \ta{வெளிப்படையாகச்} \ta{சொல்வார்}.
\end{quote}

\begin{quote}
\emph{piracciṉai iruntāl avar veḷippaṭaiyākac colvār.}
\end{quote}

\begin{quote}
If there is a problem, he will say so frankly.
\end{quote}

\ta{வெளிப்படை}, open + -\ta{ஆக}; \ta{வெளிப்படைத்தன்மை}, transparency.

\begin{cousinweb}
\textbf{\ta{வெளிப்படையாக}} (\emph{veḷippaṭaiyāka}) is native: \emph{veḷippaṭai}, \textquotedblleft{}openness\textquotedblright{}, built on \emph{veḷi}, \textquotedblleft{}outside, open\textquotedblright{} (DEDR 5498), with the adverb ending -\ta{ஆக}.
\end{cousinweb}

\begin{grammarlens}[title={What you've built}]
One more word for saying what happens if.
\end{grammarlens}

\subsection*{Your turn}
\begin{itemize}
\raggedright
  \item \emph{Say it:} \emph{veḷippaṭaiyāka}
  \item \emph{Say it:} \emph{veḷippaṭaiyāka}, once more
  \item \emph{Say it:} \emph{piracciṉai iruntāl avar veḷippaṭaiyākac colvār.} — If there is a problem, he will say so frankly.
  \item \emph{Recall:} say the Tamil for corruption, then the Tamil for a bribe, then say \emph{veḷippaṭaiyāka} again
\end{itemize}

\begin{tcolorbox}[breakable,colback=teal!4,colframe=teal!35!black,title={Before you move on}]
What does \ta{வெளிப்படையாக} mean? (\textquotedblleft{}Frankly; openly\textquotedblright{}.) Say it once more.
\end{tcolorbox}
\section[\texorpdfstring{\ta{மறைமுகமாக} (maṟaimukamāka)}{மறைமுகமாக (maṟaimukamāka)}]{\ta{மறைமுகமாக} (maṟaimukamāka) — indirectly}

\label{lesson:TA-C392-maraimukamaka}



\begin{quote}
Before the new one: say the Tamil for a bribe, then the Tamil for frankly.
\end{quote}

\subsection*{You'll want to know: \ta{மறைமுகமாக}}
\textbf{\ta{மறைமுகமாக}} — \emph{maṟaimukamāka} — \textquotedblleft{}indirectly\textquotedblright{}.

\begin{quote}
\ta{அவர்} \ta{தன்} \ta{கோபத்தை} \ta{மறைமுகமாகக்} \ta{காட்டினார்}.
\end{quote}

\begin{quote}
\emph{avar taṉ kōpattai maṟaimukamākak kāṭṭiṉār.}
\end{quote}

\begin{quote}
He showed his anger indirectly.
\end{quote}

Its opposite is \ta{நேரடியாக}, directly.

\begin{cousinweb}
\textbf{\ta{மறைமுகமாக}} (\emph{maṟaimukamāka}) is built from native \emph{maṟai}, \textquotedblleft{}to be hidden\textquotedblright{} (DEDR 4760), and \emph{mukam}, \textquotedblleft{}face\textquotedblright{}, borrowed from Sanskrit \emph{mukha} (CDIAL 10158), with -\ta{ஆக}: with a hidden face.
\end{cousinweb}

\begin{grammarlens}[title={What you've built}]
One more word for saying what happens if.
\end{grammarlens}

\subsection*{Your turn}
\begin{itemize}
\raggedright
  \item \emph{Say it:} \emph{maṟaimukamāka}
  \item \emph{Say it:} \emph{maṟaimukamāka}, once more
  \item \emph{Say it:} \emph{avar taṉ kōpattai maṟaimukamākak kāṭṭiṉār.} — He showed his anger indirectly.
  \item \emph{Recall:} say the Tamil for a bribe, then the Tamil for frankly, then say \emph{maṟaimukamāka} again
\end{itemize}

\begin{tcolorbox}[breakable,colback=teal!4,colframe=teal!35!black,title={Before you move on}]
What does \ta{மறைமுகமாக} mean? (\textquotedblleft{}Indirectly\textquotedblright{}.) Say it once more.
\end{tcolorbox}
\section[\texorpdfstring{\ta{கண்காணி} (kaṇkāṇi)}{கண்காணி (kaṇkāṇi)}]{\ta{கண்காணி} (kaṇkāṇi) — to monitor; to keep watch over}

\label{lesson:TA-C392-kankani}



\begin{quote}
Before the new one: say the Tamil for frankly, then the Tamil for indirectly.
\end{quote}

\subsection*{You'll want to know: \ta{கண்காணி}}
\textbf{\ta{கண்காணி}} — \emph{kaṇkāṇi} — \textquotedblleft{}to monitor; to keep watch over\textquotedblright{}.

\begin{quote}
\ta{தேர்தல்} \ta{அதிகாரிகள்} \ta{வாக்குச்சாவடிகளைக்} \ta{கண்காணித்தார்கள்}.
\end{quote}

\begin{quote}
\emph{tērtal atikārikaḷ vākkuccāvaṭikaḷaik kaṇkāṇittārkaḷ.}
\end{quote}

\begin{quote}
The election officials kept watch over the polling booths.
\end{quote}

Past (I …): \ta{கண்காணித்தேன்}, \emph{kaṇkāṇittēṉ}.

\ta{கண்}, eye + \ta{காணி}, the watcher: \ta{கண்காணிப்பு}, supervision.

\begin{cousinweb}
\textbf{\ta{கண்காணி}} (\emph{kaṇkāṇi}) is native: \emph{kaṇ}, \textquotedblleft{}eye\textquotedblright{} (DEDR 1159), and \emph{kāṇi}, \textquotedblleft{}one who sees\textquotedblright{}, from \emph{kāṇ}, \textquotedblleft{}to see, investigate\textquotedblright{} (DEDR 1443): to keep an eye on.
\end{cousinweb}

\begin{grammarlens}[title={What you've built}]
One more word for saying what happens if.
\end{grammarlens}

\subsection*{Your turn}
\begin{itemize}
\raggedright
  \item \emph{Say it:} \emph{kaṇkāṇi}
  \item \emph{Say it:} \emph{kaṇkāṇi}, once more
  \item \emph{Say it:} \emph{tērtal atikārikaḷ vākkuccāvaṭikaḷaik kaṇkāṇittārkaḷ.} — The election officials kept watch over the polling booths.
  \item \emph{Recall:} say the Tamil for frankly, then the Tamil for indirectly, then say \emph{kaṇkāṇi} again
\end{itemize}

\begin{tcolorbox}[breakable,colback=teal!4,colframe=teal!35!black,title={Before you move on}]
What does \ta{கண்காணி} mean? (\textquotedblleft{}To monitor; to keep watch over\textquotedblright{}.) Say it once more.
\end{tcolorbox}
\section[(dialogue)]{Looking back: from Planning Ahead to Clean Government}

\label{lesson:TA-R392-review-386-392}



\begin{quote}
Say the Tamil for a strategy and a precaution.
\end{quote}

\subsection*{Your turn}
Say these aloud:

\begin{itemize}
\raggedright
  \item a strategy, a precaution, a draft, to take forward, an aim
  \item a contract, a treaty, a promise, a pledge, to fulfil
  \item approval, to recognise officially, to appoint, to inform, a session
  \item to implement, practice, a procedure, a course of action, to administer
  \item a reform, progress, to improve, to restructure, to strengthen
  \item the economy, a subsidy, a deficit, to distribute, to fix
  \item corruption, a bribe, frankly, indirectly, to monitor
\end{itemize}

\begin{tcolorbox}[breakable,colback=teal!4,colframe=teal!35!black,title={Before you move on}]
A strategy? (\textbf{\ta{உத்தி}}, \emph{utti}.) A precaution? (\textbf{\ta{முன்னெச்சரிக்கை}}, \emph{muṉṉeccarikkai}.) A draft? (\textbf{\ta{வரைவு}}, \emph{varaivu}.) To take forward? (\textbf{\ta{முன்னெடு}}, \emph{muṉṉeṭu}.) An aim? (\textbf{\ta{குறிக்கோள்}}, \emph{kuṟikkōḷ}.) A contract? (\textbf{\ta{ஒப்பந்தம்}}, \emph{oppantam}.) A treaty? (\textbf{\ta{உடன்படிக்கை}}, \emph{uṭaṉpaṭikkai}.) A promise? (\textbf{\ta{வாக்குறுதி}}, \emph{vākkuṟuti}.) A pledge? (\textbf{\ta{உறுதிமொழி}}, \emph{uṟutimoḻi}.) To fulfil? (\textbf{\ta{நிறைவேற்று}}, \emph{niṟaivēṟṟu}.) Approval? (\textbf{\ta{ஒப்புதல்}}, \emph{opputal}.) To recognise officially? (\textbf{\ta{அங்கீகரி}}, \emph{aṅkīkari}.) To appoint? (\textbf{\ta{நியமி}}, \emph{niyami}.) To inform? (\textbf{\ta{தெரிவி}}, \emph{terivi}.) A session? (\textbf{\ta{அமர்வு}}, \emph{amarvu}.) To implement? (\textbf{\ta{செயல்படுத்து}}, \emph{ceyalpaṭuttu}.) Practice? (\textbf{\ta{நடைமுறை}}, \emph{naṭaimuṟai}.) A procedure? (\textbf{\ta{வழிமுறை}}, \emph{vaḻimuṟai}.) A course of action? (\textbf{\ta{நடவடிக்கை}}, \emph{naṭavaṭikkai}.) To administer? (\textbf{\ta{நிர்வகி}}, \emph{nirvaki}.) A reform? (\textbf{\ta{சீர்திருத்தம்}}, \emph{cīrtiruttam}.) Progress? (\textbf{\ta{முன்னேற்றம்}}, \emph{muṉṉēṟṟam}.) To improve? (\textbf{\ta{மேம்படுத்து}}, \emph{mēmpaṭuttu}.) To restructure? (\textbf{\ta{சீரமை}}, \emph{cīramai}.) To strengthen? (\textbf{\ta{வலுப்படுத்து}}, \emph{valuppaṭuttu}.) The economy? (\textbf{\ta{பொருளாதாரம்}}, \emph{poruḷātāram}.) A subsidy? (\textbf{\ta{மானியம்}}, \emph{māṉiyam}.) A deficit? (\textbf{\ta{பற்றாக்குறை}}, \emph{paṟṟākkuṟai}.) To distribute? (\textbf{\ta{விநியோகி}}, \emph{viniyōki}.) To fix? (\textbf{\ta{நிர்ணயி}}, \emph{nirṇayi}.) Corruption? (\textbf{\ta{ஊழல்}}, \emph{ūḻal}.) A bribe? (\textbf{\ta{லஞ்சம்}}, \emph{lañcam}.) Frankly? (\textbf{\ta{வெளிப்படையாக}}, \emph{veḷippaṭaiyāka}.) Indirectly? (\textbf{\ta{மறைமுகமாக}}, \emph{maṟaimukamāka}.) To monitor? (\textbf{\ta{கண்காணி}}, \emph{kaṇkāṇi}.)
\end{tcolorbox}
